\documentclass[10pt]{article}

\usepackage[utf8]{inputenc}

\usepackage[T1]{fontenc}

\usepackage{amsmath}

\usepackage{amsfonts}

\usepackage{amssymb}

\usepackage[version=4]{mhchem}

\usepackage{stmaryrd}

\usepackage{bbold}

\usepackage{graphicx}

\usepackage[export]{adjustbox}

\graphicspath{ {./images/} }



\begin{document}

$G$ - тривиальная группа. Менее тривиальный пример дает кубическая кривая $w^{2}=a_{0} z^{3}+a_{1} z^{2}+a_{2} z+a_{3}$, к-рая не допускает рациональной параметризации, но может быть униформизирована с помощью әллиптических бункций, а именно тройкой $\left(\left(f_{1}, f_{2}\right), D, G\right), \quad$ где $f_{1}$ и $f_{2}$ - рациональные функции от ф-ффункции Вейерштрасса $\mathcal{P}^{\circ}(t)$ и ее производной с соответствующими периодами $\omega_{1}, \omega_{2}$, а $G$ - группа, порожденная сдвигами $t \rightarrow t+\omega_{1}, t \rightarrow t+\omega_{2}$.



Проблема У. произвольной алгебраич. кривой, определяемой общим алгебраич. уравнением



\[

\boldsymbol{P}(z, w)=\sum_{j, k} a_{j k} z^{j} w^{k}=0,

\]



где $P$ - неприводимый алгебраич. многочлен над $\mathbb{C}$, возникла еще в 1 -й пол. 19 в., в частности в связи с интегрированием алгебраич. функций. А. Пуанкаре (H. Poincaré) поставил вопрос об У. множества репений произвольных аналитич. уравнений вида (*), когда $\boldsymbol{P}-$ сходящийся степенной ряд от двух переменных, рассматриваемый с его всевозможными аналитич. продолжениями. У. алгебраич. и произвольных аналитич. многообразий составляет содержание двадцать второй проблемы Гильберта. Полного решения проблемы У. пока (к 1984) не получено, за исключением одномерного случая.



Введя во множестве пар $(z, w)$ в $\mathbb{C}^{2}$, удовлетворяющих $(*)$, комплексную структуру с помощью элементов соответствующей алгебраич. функции $w(z)$ (или $z(w)$ ), получают компактную риманову поверхность, при этом координаты точек кривой (*) будут мероморфными функциями на этой поверхности. Более того, все компактные римановы поверхности с точностью до конформной әквивалентности получаются таким образом. Поэтому проблема У. алгебраич. кривых сводится к проблеме У. римановых поверхностей.



У. произвольной р и м а н о в о й п о в е р х н о т и $S$ наз. тройка $(D, \pi, G)$, где $D-$ область на римановой сфере $\overline{\mathbb{C}}, \pi: D \rightarrow S-$ регулярное голоморфное накрытие с накрывающей групшой $G$ конформных автоморфизмов $D$. Общая проблема заключается в нахождении и описании всех возможных таких троек для данной римановой поверхности.



Возможность У. любой римановой поверхности $S$, дающая принципиальное решение проблемы, была установлена в классич. работах П. Кебе (Р. Коеbе), А. Пуанкаре и Ф. Клейна (F. Klein); было получено окончательное решение, дающее описание всех возможных У. поверхности $S$ (см. [4] - [6]). Справедлива т е орем а унифо рмизации а нк а р е (доказанная в общем случае А. Пуанкаре, см. [2]); каждая риманова поверхность $S$ конформно эквивалентна фактору $D / G$, где $D$-одна из трех канонических областей: риманова сфера $\overline{\mathbb{C}}$, комплексная плоскость $\mathbb{C}$, единичный круг $\Delta$, а $G$ - собственно разрывная группа мёбиусовых (дробно-линейных) автоморфизмов $D$, определяемая с точностью до сопряже. ния в группе всех мёбиусовых автоморфизмов $D$.



Случаи $D=\overline{\mathbb{C}}, \mathbb{C}$ и $\Delta$ взаимно исключают друг друга. Поверхности $S$ с такими универсальными голоморфными накрытиями наз. соответственно ә л л и п т и-



\includegraphics[max width=\textwidth, center]{2023_07_09_373b4152aa7bf3b8eeabg-1}

б о л пч е с к и ми. При әтом $D=\overline{\mathbb{C}}$ только в случае, когда сама $S$ конформно эквивалентна $\overline{\mathbb{C}}$ (и, значит, $G$ тривиальна); $D=\mathbb{C}$, когда $S$ конформно әквивалентна либо $\mathbb{C}$, либо $\mathbb{C} \backslash\{0\}$, либо тору, и соответственно этому $G$ либо тривиальна, либо есть группа, порожденная сдвигом $z \rightarrow z+\omega(\omega \in \mathbb{C} \backslash\{0\})$, либо есть группа, порожденная двумя сдвигами $z \rightarrow z+\omega_{1}, z \rightarrow z+\omega_{2}$, где $\omega_{1}, \omega_{2} \neq 0$ - комілексные числа с $\operatorname{Im}\left(\omega_{2} / \omega_{1}\right) \neq$ $\neq 0$. В остальных случаях $S$ конформно әквивалентна $\Delta / G$, где $G$ - фуксова групша без кручения. Канони- ческая проекция $\pi: D \rightarrow S$ является неразветвленным накрытием и униформизирует все функции $f$ на $S$, так как функции $f \circ \pi$ однозначны в $D$. Теорема Клеӥна - Пуанкаре обобщается и на разветвленные накрытия с заданными порядками ветвления.



Несколько др. подход к проблеме У. (см. [3]) опирается на принцип: если риманова поверхность $\widetilde{S}$ гомеоморфна области $D \subset \overline{\mathbb{C}}$ (не обязательно односвязной), то $\widetilde{S}$ и конформно әквивалентна $D$. Тем самым проблема У. сводится к топологич. проблеме нахождения всех (вообще говоря, разветвленных) плоских накрытий $\tilde{S} \rightarrow S$ данной римановой поверхности $S$. Решение этих проблем дается следующими т е о р е м а м и М а с к и т a (cм. [4], [5]).



I. Пусть $S$ - ориентируемая поверхность, $v_{1}, \ldots$, $v_{n}, \ldots$ - множество простых попарно не пересекающихся петель на $S$. Если $\widetilde{S} \rightarrow S-$ регулярное накрытие с огределяющей подгрупшой $N=\left\langle v_{1}^{\alpha_{1}}, \ldots, v_{n}^{\alpha_{n}}, \ldots\right\rangle$, где $\alpha_{1}, \ldots, \alpha_{n}, .$. - натуральные числа, то $\widetilde{S}-$ плоская поверхность, т. е. гомеоморфна области в $\overline{\mathbb{C}}$.



II. Пусть $\tilde{S}-$ плоская поверхность и $\tilde{S} \rightarrow S-$ регулярное накрытие ориентируемой поверхности $S$ с определяющей подгруппой $N$. Если $S$ - поверхность конечного типа, т. е. $\pi_{1}(S)$ конечно порождена, то существуют конечное множество простых попарно не пересекающихся петель $v_{1}, \ldots, v_{n}$ и такие натуральные числа $\alpha_{1}, \ldots, \alpha_{n}$, что $\left\langle v_{1}^{\alpha_{1}}, \ldots, v_{n}^{\alpha_{n}}\right\rangle=N$.



III. Если $\tilde{S}$ - плоская риманова поверхность и $\bar{G}$ - собственно разрывная группа конформных автоморфизмов $\tilde{S}$, то существует конформный гомеоморфизм $h: \tilde{S} \rightarrow D \subset \bar{C}$ такой, что $h G h^{-1}$ есть клейнова epynna с инвариантной компонентой $D$.



Таким образом, каждая риманова поверхность униформизируется клейновой группой. Напр., если $S-$ замкнутая риманова поверхность рода $g \geqslant 1$, то ее фундаментальная группа имеет копредставление



\[

\pi_{1}(S)=\left\{a_{1}, b_{1}, \ldots, a_{g}, b_{g}: \prod_{j=1}^{g}\left[a_{j}, b_{j}\right]=1\right\}

\]



и в качестве определяющей плоское накрытие $\tilde{S}$ нормальной подгруппы $N$ можно взять наименьшую нормальную подгрупіу, порожденную $a_{1}, \ldots, a_{g}$ (или $\left.b_{1}, \ldots, b_{g}\right) ;$ тогда $S$ униформизируется $\Gamma \mathrm{p} \mathrm{у} \mathrm{ї} \mathrm{по} \mathrm{й}$ III о т к и $G$ рода $g$ - свободной чисто локсодромической клейновой группой c $g$ образующими (классич. т е о р е а $К$ ё б е о р а з р е з а $x$ ).



У. римановых поверхностей конечного типа возможные клейновы группы могут быть классифицированы. Для этого исшользуется понятие факторподгруппы. Если $G$ - клейнова группа с инвариантной компонентоӥ $D(G)$, то ее подгруппа $H$ наз. ф а к т о р по о д г р у Iгп о й $G$, если $H$ - максимальная подгрупиа, для к-рой: a) ее инвариантная компонента $D(H) \supset D(G)$ односвязна; b) $H$ не содержит с л у ч а й ны $\mathbf{x}$ п а р а б о лич е с ких э л е ме т о в $g$ (т. е. таких параболич. элементов, что при конформном изоморфизме $h: D(H) \rightarrow$ $\rightarrow \Delta$ образ $h \circ g \circ h^{-1}$ будет гиперболическим); c) каждый параболический әлемент $G$ с неподвижной точкой на предельном множестве $H$ принадлежит $H$. Напр., в теореме Клейна - Пуанкаре каждая факторподгрупша $G$ совпадает с самой $G$, а в теореме Кёбе о разрезах все факторподгруппы $G$ тривиальны. У. $(D, \pi, G)$ римановой поверхности $S$, где $D$ - инвариантная компонента $G$, наз. с т а н д а р т н о й, если $G$ не имеет кручения и не содержит случайных параболич. элементов. Для замкнутых поверхностей все такие $\mathrm{Y}$. описываются следующей теоремой (см. [6]).



Пусть $S$ - замкнутая риманова поверхность рода $g>0$ и $\left\{v_{1}, \ldots, v_{n}\right\}$ - множкество простых попарно не





\end{document}